\section[J. Time Trouble 2]{\myboxx{      \textbf{\texttt{J}}} Time Trouble 2}

\textit{Time limit}   1 second

\textit{Memory limit} 1 Gigabyte

\textit{Input} \qquad standard input

\textit{Output} \qquad standard output


\begin{quote}
	\textit{``对无数时间线建立可持久化字典树即可实现对历史的版本控制''，一位青铜龙朋友总是这么告诉我。我对此深信不疑，以至于忘记了自己倒数的心跳。}
\end{quote}

\subsection{Problem Statement}

To prove that he is never late for class, Xingyu insists that he entered the classroom at $7{:}110$ for a class beginning at $8{:}00$. This inspires a ``generalized time notation.''

In generalized time notation, a time is written $A{:}B$, where $A,B$ are nonnegative integers.

We define two \textbf{equivalence rules} for generalized times:

\begin{enumerate}
	\item For any $B \ge 60$, $A{:}B$ and $(A+1){:}(B-60)$ represent the same instant.
	\item For any $A \ge 24$, $A{:}B$ and $(A-24){:}B$ represent the same instant.
\end{enumerate}

A time $h{:}m$ is a ``standard time'' if and only if $0 \le h < 24$ and $0 \le m < 60$.

Given a generalized time, find the unique equivalent standard time.

\subsection{Input}

The first line contains an integer $T$ ($1 \le T \le 10^5$), the number of test cases.

Each test case contains two integers $A,B$ ($0 \le A,B \le 10^9$), representing the generalized time $A{:}B$.

\subsection{Output}

For each test case, output two integers $h,m$ representing the equivalent standard time $h{:}m$.

\subsection{Sample 1}

\begin{minipage}{\linewidth}
	\begin{minipage}[t]{0.49\linewidth}
		\begin{lstlisting}[title=\texttt{standard input}]
5
7 110
25 30
0 0
23 59
100 200
\end{lstlisting}
	\end{minipage}
	\begin{minipage}[t]{0.49\linewidth}
		\begin{lstlisting}[title=\texttt{standard output}]
8 50
1 30
0 0
23 59
7 20
\end{lstlisting}
	\end{minipage}
\end{minipage}


\subsection{Sample 1 Explanation}

In the first test case, $7{:}110 \to 8{:}50$ by rule 1, since $110=60+50$.

In the second test case, $25{:}30 \to 1{:}30$ by rule 2, since $25=24+1$.

In the third test case, $0{:}0$ is already a standard time.

